% ==========================================================
%  OLD EXAM QUESTIONS — DL Architectures & Generative Techniques
%  (WORKED SOLUTIONS)
%  Companion to exams.tex: every question is followed
%  immediately by a green answer box. For multiple-choice,
%  correct options are marked with a green check, incorrect
%  ones with a red cross, each with a short justification
%  (incl. why tricky near-miss options are wrong). For
%  calculations the working is shown; wild distractors that
%  only exist to catch guessers get no explanation.
%
%  NOTE: exams.tex stays the blank practice set. This file is
%  the solutions variant — do not merge the two.
% ==========================================================

\documentclass[11pt,oneside]{article}

\usepackage[
  a4paper,
  top=2.4cm,
  bottom=2.1cm,
  left=2.3cm,
  right=2.3cm,
  headheight=15pt
]{geometry}

\usepackage[T1]{fontenc}
\usepackage[utf8]{inputenc}
\usepackage{lmodern}
\usepackage{microtype}
\usepackage{inconsolata}
\usepackage{amsmath,amssymb,mathtools,bm}
\usepackage{enumitem}
\usepackage{booktabs,array,multicol,longtable}
\usepackage{xcolor}
\usepackage{graphicx}
\usepackage[most]{tcolorbox}
\usepackage{titlesec}
\usepackage{fancyhdr}
\usepackage[hidelinks]{hyperref}

\setlength{\parindent}{0pt}
\setlength{\parskip}{0.55em}
\setlength{\arraycolsep}{2.5pt}   % tighten matrices in answer options
\setlist[itemize]{topsep=3pt,itemsep=2pt,leftmargin=2.2em}
\setlist[enumerate]{topsep=3pt,itemsep=2pt,leftmargin=2.2em}

% ── Colors ───────────────────────────────────────────────
% Shared accent with summary.tex — do not change.
\definecolor{accent}{HTML}{1B3A5C}
\definecolor{q@frame}{HTML}{1A5276}
\definecolor{q@bg}{HTML}{F6FAFD}
\definecolor{key@frame}{HTML}{1E8449}
\definecolor{key@bg}{HTML}{F4FBF6}
\definecolor{note@frame}{HTML}{9A6A00}
\definecolor{note@bg}{HTML}{FEF9E7}
\definecolor{ok}{HTML}{1E8449}    % green check — correct option
\definecolor{bad}{HTML}{B03A2E}   % red cross  — incorrect option

% ── Box global style ─────────────────────────────────────
\tcbset{
  enhanced,
  breakable,
  arc=2.5mm,
  boxrule=0.5pt,
  left=4.5mm,
  right=4mm,
  top=2.5mm,
  bottom=2.5mm,
  fonttitle=\bfseries\small,
  sharp corners=west,
  parbox=false
}

% Question box (blue) — one per exam question.
\newtcolorbox{qbox}[1][]{
  colback=q@bg,
  colframe=q@frame,
  borderline west={3.5pt}{0pt}{q@frame},
  title={#1}
}

% Key insight box (green) — answer hints, key results.
\newtcolorbox{keybox}[1][]{
  colback=key@bg,
  colframe=key@frame,
  borderline west={3.5pt}{0pt}{key@frame},
  title={#1}
}

% Organizational note box (amber) — structural remarks only.
\newtcolorbox{notebox}[1][]{
  colback=note@bg,
  colframe=note@frame,
  borderline west={3.5pt}{0pt}{note@frame},
  title={#1}
}

% ── Section headings ─────────────────────────────────────
\titleformat{\section}
  {\normalfont\Large\bfseries\color{accent}}
  {}{0em}{}
  [{\vspace{1pt}\color{accent!45}\titlerule[0.45pt]}]
\titlespacing*{\section}{0pt}{3.2ex plus 1ex minus .2ex}{1.8ex plus .2ex}

\titleformat{\subsection}
  {\normalfont\normalsize\bfseries\color{accent!80!black}}
  {}{0em}{}
\titlespacing*{\subsection}{0pt}{2.2ex plus 1ex minus .2ex}{1.0ex plus .2ex}

% ── Header & footer ──────────────────────────────────────
\pagestyle{fancy}
\fancyhf{}
\fancyhead[L]{\small\color{gray!70}Old Exam Questions (Solutions) — \coursetitle}
\fancyhead[R]{\small\color{gray!70}\thepage}
\renewcommand{\headrulewidth}{0.4pt}

% ── Document metadata ─────────────────────────────────────
\newcommand{\coursetitle}{DL Architectures \& Generative Techniques}

% ── Question counter & macros ────────────────────────────
\newcounter{question}
\newcommand{\questiontitle}[1]{%
  \stepcounter{question}\begin{qbox}[Question \thequestion: #1]}
\newcommand{\closequestion}{\end{qbox}}

% Mark questions that appeared on multiple exams.
\newcommand{\repeatnote}{\textit{\color{gray!70}Asked multiple times.}}

% Multiple-choice checkbox bullet.
\newcommand{\choice}{\item[$\square$] }

\newcommand{\examfigure}[2][width=0.82\linewidth]{%
  \begin{center}
    \includegraphics[#1]{#2}
  \end{center}
}

% ── Inline answer macros (solutions variant) ─────────────
% Green answer box placed directly under each question.
% Multiple-choice usage:
%   \begin{answer}{A, B}
%     \yes{A} reason why correct ...
%     \no{C}  reason why wrong (esp. for tricky near-misses) ...
%   \end{answer}
\newenvironment{answer}[1]{%
  \begin{keybox}[Answer]%
  \begin{itemize}[leftmargin=1.7em,itemsep=1.5pt,topsep=2pt]}%
  {\end{itemize}\end{keybox}}
\newcommand{\yes}[1]{\item[\textcolor{ok}{\ensuremath{\checkmark}}]\textbf{#1.}\ }
\newcommand{\no}[1]{\item[\textcolor{bad}{\ensuremath{\bm\times}}]\textbf{#1.}\ }

% Prose answer box — for calculations and labelling questions.
% Usage:  \begin{answerprose}{result}  working / explanation  \end{answerprose}
\newenvironment{answerprose}[1]{%
  \begin{keybox}[Answer]}%
  {\end{keybox}}

% ── Math macros ──────────────────────────────────────────
\newcommand{\R}{\mathbb{R}}
\newcommand{\E}{\mathbb{E}}
\newcommand{\Prob}{\mathbb{P}}
\newcommand{\KL}[2]{D_{\mathrm{KL}}\!\left(#1\,\|\,#2\right)}
\newcommand{\dd}{\,\mathrm{d}}
\newcommand{\Cov}{\mathrm{Cov}}
\newcommand{\loss}{\mathcal{L}}
\newcommand{\Normal}{\mathcal{N}}
\newcommand{\given}{\,\vert\,}
\newcommand{\T}{^{\mathsf{T}}}
\DeclareMathOperator{\softmax}{softmax}


% ==========================================================
%  BEGIN DOCUMENT
% ==========================================================
\begin{document}

% ── Title page ───────────────────────────────────────────
\thispagestyle{empty}
\vspace*{1.4cm}
\begin{center}
  {\color{accent}\rule{\linewidth}{1.8pt}}\par
  \vspace{0.65cm}
  {\fontsize{26}{32}\selectfont\bfseries\color{accent}Old Exam Questions\\[2pt] \coursetitle\par}
  \vspace{0.25cm}
  {\Large\color{key@frame}\bfseries Worked Solutions\par}
  \vspace{0.55cm}
  {\color{accent!35}\rule{\linewidth}{0.5pt}}\par
  \vspace{0.55cm}
  {\large Stefan Eder\par}
  \vspace{0.15cm}
  {\small\color{gray!70}\today\par}
\end{center}

\vspace{0.9cm}
\begin{notebox}[How this file is organized]
This is the \textbf{solutions companion} to \texttt{exams.tex}, collected from the main exam
(28.06.2024) and the 2022 / 2024 retry exams. Every question is followed \emph{immediately} by a
green \textbf{answer box}. For \textbf{multiple-choice} questions (``select one or more'') each
option is marked \textcolor{ok}{\ensuremath{\checkmark}} (correct) or
\textcolor{bad}{\ensuremath{\bm\times}} (incorrect), with a short justification --- \emph{why} the
correct options are correct and, for the tricky / near-miss distractors, \emph{why} they are wrong.
For \textbf{calculations} the working is shown; distractor options that exist only to catch a guess
get no explanation. If you want a clean practice set with blank answers, use \texttt{exams.tex}.
\end{notebox}

\tableofcontents
\newpage


% ==========================================================
\section{Image Classification \& CNN Architectures}
% ==========================================================

\questiontitle{AlexNet}
Which of the following statements about AlexNet (Krizhevsky et al., 2012) are correct?
(Select one or more.)
\begin{itemize}[leftmargin=1.8em]
  \choice AlexNet does not have fully-connected layers.
  \choice AlexNet uses strided convolutions instead of max-pooling for down-sampling.
  \choice AlexNet is a deep convolutional network with more than 1 hidden layer.
  \choice One of the most important components of AlexNet is BatchNorm.
  \choice AlexNet uses input images with a resolution of $224\times224$ pixels and three channels.
\end{itemize}
\closequestion
\begin{answer}{C, E}
  \no{A} AlexNet ends in three large \emph{fully-connected} layers --- it definitely has them.
  \no{B} It down-samples with \emph{max-pooling} (alongside convolutions); it does not replace pooling with strided convs.
  \yes{C} It is a deep CNN --- 5 convolutional + 3 fully-connected layers.
  \no{D} BatchNorm did not exist yet (2015). AlexNet's key novelties were ReLU, dropout and local response normalization --- not BatchNorm.
  \yes{E} Input is $224\times224\times3$ RGB.
\end{answer}

\questiontitle{VGGNet}
Which of the following statements about VGGNet are correct? (Select one or more.)
\begin{itemize}[leftmargin=1.8em]
  \choice It predominantly uses $3\times3$ convolutions.
  \choice It is a single-stage approach (a single $5\times5$ layer rather than stacked $3\times3$ layers).
  \choice It has a symmetric encoder--decoder structure.
  \choice The majority of the network's weights are located in the fully-connected layers.
  \choice It cannot be used for image classification.
\end{itemize}
\closequestion
\begin{answer}{A, D}
  \yes{A} VGG's signature is deep stacks of small $3\times3$ convolutions.
  \no{B} The opposite: VGG \emph{stacks} two $3\times3$ layers to match one $5\times5$'s receptive field with fewer parameters and an extra non-linearity.
  \no{C} It is a plain feed-forward classifier, \emph{not} an encoder--decoder.
  \yes{D} The huge FC layers (especially the first) hold the bulk of the $\sim$138M parameters.
  \no{E} Nonsense --- VGG is an ImageNet classifier.
\end{answer}

\questiontitle{Normalization in Image Classifiers}
Which of the following statements about normalization in image-classification networks are correct?
(Select one or more.)
\begin{itemize}[leftmargin=1.8em]
  \choice Batch Normalization has been used since AlexNet.
  \choice Since Xception, most image classifiers started to use Layer Normalization.
  \choice Batch Normalization was used in InceptionNet to speed up training.
  \choice Batch Normalization is part of most deep networks for image classification (e.g.\ ResNet).
\end{itemize}
\closequestion
\begin{answer}{C, D}
  \no{A} AlexNet used \emph{Local Response Normalization}; BatchNorm only appeared in 2015.
  \no{B} LayerNorm is mainly a transformer / sequence-model thing --- CNN classifiers still use BatchNorm.
  \yes{C} Inception / GoogLeNet introduced BatchNorm precisely to accelerate training.
  \yes{D} BatchNorm is a standard component of deep CNN classifiers such as ResNet.
\end{answer}

\questiontitle{Residual Networks}
Which of the following statements about Residual Networks (ResNets) are correct? (Select one or more.)
\begin{itemize}[leftmargin=1.8em]
  \choice Residual networks use skip-connections that add the representation $\mathbf{h}^{l}$ of a
        former layer $l$ to a later layer $l+1$, i.e.\ $\mathbf{h}^{l+1}=\mathbf{h}^{l}+F(\mathbf{h}^{l})$,
        where $F$ can be any transformation with appropriate dimensions (e.g.\ a fully-connected layer).
  \choice Because of skip-connections, very deep ResNets can be trained even without normalization techniques.
  \choice Pre-trained ResNets cannot be used for transfer learning because the learned features are usually
        overly specified to the task on which they were trained.
  \choice Residual networks are always convolutional and cannot be used in purely fully-connected networks.
  \choice Leaky ReLUs cannot be used in Residual Networks.
\end{itemize}
\closequestion
\begin{answer}{A}
  \yes{A} That is exactly the residual formulation; $F$ can be any dimension-matching transform (conv, FC, $\dots$).
  \no{B} Normalization is still central to training very deep ResNets --- skip-connections alone don't remove that need.
  \no{C} Pre-trained ResNets are one of the most common transfer-learning backbones; their features are general, not over-specialised.
  \no{D} Residual connections work in fully-connected nets too --- they are not convolution-only.
  \no{E} Any activation, leaky ReLU included, can be used inside a residual block.
\end{answer}

\questiontitle{$1\times1$ Convolution Output Size}
You add a $1\times1$ convolutional layer with 16 kernels, padding 0, stride 1, to an input of
shape $[128,128,64]$. Which statement is correct?
\begin{itemize}[leftmargin=1.8em]
  \choice It is not possible to apply this convolution layer.
  \choice The resulting dimension is $[128,128,16]$.
  \choice A $1\times1$ convolution does not exist.
  \choice Better: replace the $1\times1$ layer with $2\times2$ max-pooling --- an equivalent but more efficient operation.
  \choice Better: use $1\times1$ with stride 2 --- equivalent but more efficient.
\end{itemize}
\closequestion
\begin{answer}{B}
  \no{A} It is perfectly valid.
  \yes{B} A $1\times1$ conv keeps the spatial size ($128\times128$) and sets the channel depth to the number of kernels (16): $[128,128,16]$.
  \no{C} Nonsense --- $1\times1$ convolutions are standard (channel mixing / bottlenecks).
  \no{D} $2\times2$ max-pooling halves the spatial resolution --- not equivalent.
  \no{E} Stride 2 would down-sample to $64\times64$ --- not equivalent.
\end{answer}

\questiontitle{Receptive Field of the Last Layer}
You are given a CNN with a fixed architecture. You may only change the parameters of the
\emph{last} convolutional layer, and you want to strongly increase the receptive field (RF)
of each neuron in that last layer. For each change, decide whether it produces
\textbf{no increase}, a \textbf{small increase}, or a \textbf{strong increase} of the RF.
\begin{itemize}[leftmargin=1.8em]
  \item Increase the stride of the convolution
  \item Switch to dilated convolutions
  \item Increase the number of kernels
  \item Increase the kernel size
\end{itemize}
\closequestion
\begin{answerprose}{no / strong / no / small}
\begin{itemize}[leftmargin=1.7em,itemsep=1.5pt,topsep=2pt]
  \item[\textcolor{ok}{$\to$}] \textbf{Increase the stride:} \emph{no increase.} Stride in the \emph{last} layer only sub-samples its output; it does not enlarge the input region each of those neurons already sees.
  \item[\textcolor{ok}{$\to$}] \textbf{Dilated convolutions:} \emph{strong increase.} Dilation spreads the kernel taps far apart, sharply enlarging the RF.
  \item[\textcolor{ok}{$\to$}] \textbf{Number of kernels:} \emph{no increase.} More channels means more features, identical spatial RF.
  \item[\textcolor{ok}{$\to$}] \textbf{Kernel size:} \emph{small increase.} A larger kernel adds a few pixels of RF, but only in this one layer.
\end{itemize}
\end{answerprose}

\questiontitle{Zero Initialization --- Forward Pass}
You design a standard 2-layer fully-connected network with sigmoid activations and SGD.
You initialize all weights and biases to zero and forward-propagate an input $\mathbf{x}$.
What is the output $\hat y$?
\begin{itemize}[leftmargin=1.8em]
  \choice $0$
  \choice $e^{x}$
  \choice $0.5$
  \choice $-1$
  \choice $1$
\end{itemize}
\closequestion
\begin{answer}{C}
  \no{A} Not $\sigma(0)$.
  \no{B} Nonsense.
  \yes{C} With all weights and biases zero, every pre-activation is $0$, and $\sigma(0)=0.5$ at every layer --- so $\hat y = 0.5$.
  \no{D} Out of range for a sigmoid.
  \no{E} Not $\sigma(0)$.
\end{answer}

\questiontitle{Zero Initialization --- Weight Update}
Consider the same all-zero-initialized network. You forward-propagate a batch, backpropagate,
and update the parameters once. $\mathbf{W}^{(2)}$ denotes the weight matrix of the second layer.
Which statement is then true?
\begin{itemize}[leftmargin=1.8em]
  \choice The entries of $\mathbf{W}^{(2)}$ are all positive.
  \choice The entries of $\mathbf{W}^{(2)}$ may be positive or negative.
  \choice The entries of $\mathbf{W}^{(2)}$ are all negative.
  \choice The entries of $\mathbf{W}^{(2)}$ are all zero.
\end{itemize}
\closequestion
\begin{answer}{B}
  \no{A} The sign is not forced positive.
  \yes{B} The hidden activations are all the same positive value ($0.5$), but the update sign is set by the per-output error $(\hat y - y)$, which can be positive or negative.
  \no{C} The sign is not forced negative.
  \no{D} The near-miss: the second-layer weights \emph{do} get a non-zero gradient because their input ($0.5$) is non-zero --- so they don't stay at zero. (The \emph{first}-layer weights would still be stuck.)
\end{answer}


% ==========================================================
\section{Semantic Segmentation}
% ==========================================================

\questiontitle{Self-Driving Scene Labelling}
A camera in a self-driving car provides road-scene images, and you train a network that
outputs an image in which every pixel is annotated as car, road, pedestrian, etc.
The figure shows examples of the input images and the corresponding pixel-wise labels.
\examfigure[width=0.88\linewidth]{figures/exam-semantic-segmentation-road-scene.png}
Which statements are true in this context? (Select one or more.)
\begin{itemize}[leftmargin=1.8em]
  \choice To train such a network, labels for each pixel of the training images are required.
  \choice This represents an object-detection task in computer vision.
  \choice Generative adversarial networks are the most appropriate type of network for this task.
  \choice Deep networks cannot be used because there is no appropriate loss function.
  \choice The intersection over union (IoU) would not be an appropriate metric to assess such a model.
\end{itemize}
\closequestion
\begin{answer}{A}
  \yes{A} Per-pixel labels are needed --- this is \emph{semantic segmentation}.
  \no{B} It assigns a class to every pixel, not bounding boxes --- not object detection.
  \no{C} A CNN (e.g.\ U-Net / FCN) with a per-pixel cross-entropy loss is the natural choice, not a GAN.
  \no{D} There is a perfectly good loss (per-pixel cross-entropy) --- nonsense.
  \no{E} IoU / mIoU is \emph{the} standard segmentation metric, so it is appropriate.
\end{answer}

\questiontitle{IoU Computation}
The figure shows a ground-truth object region (solid brown outline) and a predicted object region
(dashed green outline) on a grid of pixels. The ground-truth region covers $10$ pixels, the
predicted region covers $8$ pixels, and they overlap in $5$ pixels. Compute the IoU metric,
rounded to three decimals.
\examfigure[width=0.58\linewidth]{figures/exam-iou-pixel-regions.png}
\closequestion
\begin{answerprose}{$5/13 \approx 0.385$}
$\mathrm{IoU}=\dfrac{|\text{intersection}|}{|\text{union}|}=\dfrac{5}{10+8-5}=\dfrac{5}{13}\approx 0.385$.
The union subtracts the double-counted overlap.
\end{answerprose}


% ==========================================================
\section{Object Detection \& Localization}
% ==========================================================

\questiontitle{Object-Detection Methods}
Which of the following statements about deep-learning-based object-detection methods are correct?
(Select one or more.)
\begin{itemize}[leftmargin=1.8em]
  \choice Faster R-CNN uses a region proposal network rather than the selective-search algorithm used in R-CNN.
  \choice The operations RoI-pooling and RoI-align operate on feature maps and not on the input image.
  \choice R-CNN uses AlexNet to extract features for Regions of Interest.
  \choice Object detection can be formulated as a regression task in a multi-task setting.
  \choice Fast R-CNN has two output layers (heads): one for regression and one for classification.
\end{itemize}
\closequestion
\begin{answer}{A, B, D, E}
  \yes{A} Faster R-CNN's contribution is the learned \emph{Region Proposal Network}, replacing selective search.
  \yes{B} Both RoI-pooling and RoI-align pool from the convolutional \emph{feature maps}, not the raw image.
  \no{C} R-CNN does run a CNN on each warped proposal, but it is not characterised as ``AlexNet'' in the lecture --- treated as false here.
  \yes{D} Detection is a multi-task problem: box-coordinate \emph{regression} + class prediction.
  \yes{E} Fast R-CNN has a bounding-box regression head and a classification head.
\end{answer}

\questiontitle{One-Stage vs.\ Two-Stage Detectors}
For each architecture, state whether it is a \textbf{one-stage} or \textbf{two-stage} detector.
\begin{itemize}[leftmargin=1.8em]
  \item Faster R-CNN
  \item YOLO
  \item Mask R-CNN
  \item Fast R-CNN
  \item R-CNN
  \item SSD
\end{itemize}
\closequestion
\begin{answerprose}{the R-CNN family is two-stage; YOLO and SSD are one-stage}
\begin{itemize}[leftmargin=1.7em,itemsep=1.5pt,topsep=2pt]
  \item[\textcolor{ok}{$\to$}] Faster R-CNN --- \textbf{two}-stage \quad (propose regions, then classify/refine)
  \item[\textcolor{ok}{$\to$}] YOLO --- \textbf{one}-stage
  \item[\textcolor{ok}{$\to$}] Mask R-CNN --- \textbf{two}-stage
  \item[\textcolor{ok}{$\to$}] Fast R-CNN --- \textbf{two}-stage
  \item[\textcolor{ok}{$\to$}] R-CNN --- \textbf{two}-stage
  \item[\textcolor{ok}{$\to$}] SSD --- \textbf{one}-stage
\end{itemize}
Rule of thumb: anything in the R-CNN family has a separate proposal stage (two-stage); YOLO/SSD predict boxes and classes in a single pass (one-stage).
\end{answerprose}

\questiontitle{Transpose Convolution (stride 1)}
You are given an input $\mathbf{x}=\begin{pmatrix}0&1\\2&3\end{pmatrix}$ and a kernel
$\mathbf{w}=\begin{pmatrix}0&1\\2&3\end{pmatrix}$. What is the result of the transpose
convolution $\mathbf{w}\ast^{T}\mathbf{x}$ (stride 1)?
{\footnotesize
\begin{itemize}[leftmargin=1.8em]
  \choice $(42)$
  \choice $(14)$
  \choice $\begin{pmatrix}1&1&1&1\\2&2&2&2\\3&3&3&3\\4&4&4&4\end{pmatrix}$
  \choice $\begin{pmatrix}13&13&13\\13&13&13\\13&13&13\end{pmatrix}$
  \choice $\begin{pmatrix}0&2&1\\6&14&6\\2&3&0\end{pmatrix}$
  \choice $\begin{pmatrix}3&2\\1&0\end{pmatrix}$
  \choice $\begin{pmatrix}0&0&0&1\\0&0&2&3\\0&2&0&3\\4&6&6&9\end{pmatrix}$
  \choice $\begin{pmatrix}0&0&1\\0&4&6\\4&12&9\end{pmatrix}$
\end{itemize}}
\closequestion
\begin{answerprose}{H, $\begin{psmallmatrix}0&0&1\\0&4&6\\4&12&9\end{psmallmatrix}$}
A stride-1 transpose convolution scatters a copy of the kernel, scaled by each input entry, and
overlap-adds them. Output size $(2-1)\cdot 1 + 2 = 3$, i.e.\ $3\times3$:
\[
1\!\cdot\!\begin{psmallmatrix}0&1\\2&3\end{psmallmatrix}\!\text{ at }(0,1)
\;+\;2\!\cdot\!\begin{psmallmatrix}0&1\\2&3\end{psmallmatrix}\!\text{ at }(1,0)
\;+\;3\!\cdot\!\begin{psmallmatrix}0&1\\2&3\end{psmallmatrix}\!\text{ at }(1,1)
=\begin{psmallmatrix}0&0&1\\0&4&6\\4&12&9\end{psmallmatrix}.
\]
The other options are wrong shapes / values that only test whether you actually carried out the overlap-add.
\end{answerprose}

\questiontitle{Transpose Convolution (stride 2)}
You are given an input $\mathbf{x}=\begin{pmatrix}1&-2\\-1&2\end{pmatrix}$ and a kernel
$\mathbf{w}=\begin{pmatrix}0&1\\2&3\end{pmatrix}$. What is the result of the transpose
convolution $\mathbf{w}\ast^{T}\mathbf{x}$ with \textbf{stride 2}?
{\footnotesize
\begin{itemize}[leftmargin=1.8em]
  \choice $\begin{pmatrix}0&1&-1\\2&-1&-1\\-4&-2&6\end{pmatrix}$
  \choice $\begin{pmatrix}3&2\\1&0\end{pmatrix}$
  \choice $\begin{pmatrix}0&1&0&-2\\2&3&-4&-6\\0&-1&0&2\\-2&-3&4&6\end{pmatrix}$
  \choice $\begin{pmatrix}0&2&1\\6&-14&6\\2&3&0\end{pmatrix}$
  \choice $(3)$
  \choice $(42)$
\end{itemize}}
\closequestion
\begin{answerprose}{C, $\begin{psmallmatrix}0&1&0&-2\\2&3&-4&-6\\0&-1&0&2\\-2&-3&4&6\end{psmallmatrix}$}
With stride 2, the (size-2) kernel copies are placed two apart, so they \emph{don't overlap}; output size
$(2-1)\cdot 2 + 2 = 4$, i.e.\ $4\times4$. Each block is $x_{ij}\cdot\mathbf{w}$ at offset $(2i,2j)$:
\[
\begin{array}{c|c}
1\!\cdot\!\begin{psmallmatrix}0&1\\2&3\end{psmallmatrix} & -2\!\cdot\!\begin{psmallmatrix}0&1\\2&3\end{psmallmatrix}\\[6pt]\hline
-1\!\cdot\!\begin{psmallmatrix}0&1\\2&3\end{psmallmatrix} & 2\!\cdot\!\begin{psmallmatrix}0&1\\2&3\end{psmallmatrix}
\end{array}
=\begin{psmallmatrix}0&1&0&-2\\2&3&-4&-6\\0&-1&0&2\\-2&-3&4&6\end{psmallmatrix}.
\]
The remaining options are wrong sizes / values to catch a guess.
\end{answerprose}


% ==========================================================
\section{Autoencoders}
% ==========================================================

\questiontitle{Annotating the Generative-Model Taxonomy}
A figure shows three categories of generative models. Each row is a network with grey input/output
boxes and a latent variable $\mathbf{z}$ in the middle. Annotate the building blocks with
\emph{Generator}, \emph{Discriminator}, \emph{Encoder}, \emph{Decoder}, \emph{Flow}, \emph{Inverse Flow},
given that a rectangular block keeps the dimension constant while a trapezoidal block decreases or
increases it.
\examfigure[width=0.9\linewidth]{figures/exam-generative-model-taxonomy.png}
\begin{itemize}[leftmargin=1.8em]
  \item Top row: input $\to$ trapezoid $\to 0/1 \to \mathbf{z} \to$ trapezoid $\to$ output
  \item Middle row: input $\to$ trapezoid (down) $\to \mathbf{z} \to$ trapezoid (up) $\to$ output
  \item Bottom row: input $\to$ rectangle $\to \mathbf{z} \to$ rectangle $\to$ output
\end{itemize}
\closequestion
\begin{answerprose}{GAN / (V)AE / Flow}
\begin{itemize}[leftmargin=1.7em,itemsep=1.5pt,topsep=2pt]
  \item[\textcolor{ok}{$\to$}] \textbf{Top row = GAN.} The block ending in a $0/1$ decision is the \textbf{Discriminator}; the block going from $\mathbf{z}$ to the output is the \textbf{Generator}.
  \item[\textcolor{ok}{$\to$}] \textbf{Middle row = (V)AE.} The down-sizing trapezoid is the \textbf{Encoder} $q(\mathbf{z}\given\mathbf{x})$; the up-sizing one is the \textbf{Decoder} $p(\mathbf{x}\given\mathbf{z})$.
  \item[\textcolor{ok}{$\to$}] \textbf{Bottom row = normalizing flow.} Dimensions are preserved (rectangles): the first block is the \textbf{Flow}, the second its \textbf{Inverse Flow}.
\end{itemize}
\end{answerprose}

\questiontitle{Autoencoders (excluding VAEs)}
Which of the following statements about autoencoders (\textbf{not} including variational autoencoders)
are correct? (Select one or more.)
\begin{itemize}[leftmargin=1.8em]
  \choice The input--output Jacobian of an autoencoder is always a lower-triangular matrix.
  \choice Autoencoders do not map a single input to a single point in latent space, but to a distribution
        in latent space.
  \choice Autoencoders are always fully-connected networks.
  \choice All autoencoders use fewer neurons in the hidden layers than in the input layer.
  \choice Autoencoders consist of an encoder and a decoder, which are parametrized neural networks.
\end{itemize}
\closequestion
\begin{answer}{E}
  \no{A} There is no such constraint --- the Jacobian is generally not triangular.
  \no{B} Mapping an input to a \emph{distribution} is the VAE; a plain AE maps it to a single latent \emph{point}.
  \no{C} They can be convolutional, recurrent, etc.\ --- not necessarily fully-connected.
  \no{D} Undercomplete is common but not required: overcomplete / sparse AEs have $\ge$ input-size hidden layers.
  \yes{E} The only always-true statement: an AE is an encoder + decoder, both parametrized networks.
\end{answer}

\questiontitle{Sparse and Denoising Autoencoders}
Which of the following statements about sparse / denoising autoencoders are correct?
(Select one or more.)
\begin{itemize}[leftmargin=1.8em]
  \choice Sparse AEs enforce sparsity on the representation via a regularization term.
  \choice Sparse AEs enforce hidden units to be inactive (i.e.\ activation $=0$) for most inputs.
  \choice A sparse AE can be trained on a copy task, where the input equals the target.
  \choice Sparse AEs force the \emph{input} units to be inactive (activation $=0$).
  \choice Sparse AEs enforce sparsity via early stopping.
  \choice A sparse AE cannot be trained on a denoising task (input noisy, target clean).
\end{itemize}
\closequestion
\begin{answer}{A, B, C}
  \yes{A} Sparsity is added as a regularization term (e.g.\ a KL penalty on mean activation, or $L_1$).
  \yes{B} The effect is that most \emph{hidden} units are (near-)zero for any given input.
  \yes{C} It can be trained on the ordinary copy / reconstruction task (input $=$ target).
  \no{D} It silences \emph{hidden} units, not input units --- the input is the data.
  \no{E} Sparsity comes from the penalty term, not from early stopping.
  \no{F} Nothing stops combining a sparsity penalty with a denoising objective.
\end{answer}

\questiontitle{Autoencoder Reconstruction Loss}
An autoencoder uses ReLU activations with encoder
$u=\mathrm{ReLU}\!\left(\begin{bmatrix}1&0&-1\\0&1&0\end{bmatrix}\mathbf{x}\right)$ and decoder
$\hat{\mathbf{x}}=\mathrm{ReLU}\!\left(\begin{bmatrix}1&0\\0&0\\1&0\end{bmatrix}u\right)$.
Compute the reconstruction loss $\loss(\mathbf{x},\hat{\mathbf{x}})=\sum_i (x_i-\hat x_i)^2$
for the input $\mathbf{x}=(1,1,2)\T$.
\closequestion
\begin{answerprose}{$6$}
Encoder: $\begin{bmatrix}1&0&-1\\0&1&0\end{bmatrix}(1,1,2)\T=(1-2,\,1)=(-1,1)$, so
$u=\mathrm{ReLU}(-1,1)=(0,1)$.\\
Decoder: $\begin{bmatrix}1&0\\0&0\\1&0\end{bmatrix}(0,1)\T=(0,0,0)$, so $\hat{\mathbf{x}}=\mathrm{ReLU}(0,0,0)=(0,0,0)$.\\
Loss $=(1-0)^2+(1-0)^2+(2-0)^2=1+1+4=6$. \quad (If a $\tfrac1n$ MSE were used instead of the sum, the value would be $2.0$.)
\end{answerprose}


% ==========================================================
\section{Variational Autoencoders}
% ==========================================================

\questiontitle{Variational Autoencoders}
Which of the following statements about variational autoencoders (VAEs) are true? (Select one or more.)
\begin{itemize}[leftmargin=1.8em]
  \choice VAEs have an encoder--decoder structure, where encoder and decoder must be symmetric.
  \choice The prior distribution of the latent variables $\mathbf{z}$ must be a discrete distribution.
  \choice VAEs maximize a lower bound on the likelihood of the data.
  \choice VAEs automatically provide the marginal likelihood $p(\mathbf{x})$ at low computational cost.
  \choice VAEs cannot be used to generate images.
\end{itemize}
\closequestion
\begin{answer}{C}
  \no{A} Encoder and decoder need not be symmetric.
  \no{B} The prior is typically \emph{continuous} (a standard Gaussian), not discrete.
  \yes{C} VAEs maximize the ELBO, a lower bound on the data log-likelihood.
  \no{D} They do \emph{not} give $p(\mathbf{x})$ cheaply --- the marginal stays intractable; the ELBO is only a bound.
  \no{E} Nonsense --- generating images is a primary use of VAEs.
\end{answer}

\questiontitle{VAE Objective Function}
Using the notation from the lecture, which of the following is the correct formulation of the
standard Variational Autoencoder (VAE) objective?
\begin{itemize}[leftmargin=1.8em]
  \choice $\loss(D,G)=\E_{\mathbf{x}\sim p_r}[\log D(\mathbf{x})]+\E_{\mathbf{z}\sim p_z}[\log(1-D(G(\mathbf{z})))]$
  \choice $\loss_{\theta,\phi}(\mathbf{x})=\E_{q_\phi(\mathbf{z}\given\mathbf{x})}\log p_\theta(\mathbf{x}\given\mathbf{z})-\KL{p_\theta(\mathbf{x})}{p_\theta(\mathbf{z})}$
  \choice $\loss(D,G)=\E_{\mathbf{x}\sim p_r}[\log D(\mathbf{x})]+\E_{\mathbf{z}\sim p_z}[D(G(\mathbf{z}))]$
  \choice $\loss_{\theta,\phi}(\mathbf{x})=\E_{q_\phi(\mathbf{z}\given\mathbf{x})}\log p_\theta(\mathbf{x}\given\mathbf{z})-\KL{q_\phi(\mathbf{z}\given\mathbf{x})}{p_\theta(\mathbf{z})}$
  \choice $\loss_\theta(\mathbf{x})=\KL{p_{\mathrm{data}}(\mathbf{x};\theta)}{p^{*}(\mathbf{x})}$
\end{itemize}
\closequestion
\begin{answer}{D}
  \no{A} This is the GAN min--max value function, not the VAE objective.
  \no{B} Right shape, \emph{wrong KL arguments}: the KL must be between $q_\phi(\mathbf{z}\given\mathbf{x})$ and the prior $p_\theta(\mathbf{z})$, not between $p_\theta(\mathbf{x})$ and $p_\theta(\mathbf{z})$.
  \no{C} Another GAN-style objective (and missing the $\log$ on the second term).
  \yes{D} The ELBO: reconstruction term $\E_{q_\phi(\mathbf{z}\given\mathbf{x})}\log p_\theta(\mathbf{x}\given\mathbf{z})$ minus $\KL{q_\phi(\mathbf{z}\given\mathbf{x})}{p_\theta(\mathbf{z})}$.
  \no{E} Nonsense --- not the VAE objective.
\end{answer}

\questiontitle{Best Gaussian Approximation via KL}
\repeatnote\\
You have a fixed, multi-dimensional, non-Gaussian distribution $p(\mathbf{x})$ that you want to
approximate with a Gaussian $q(\mathbf{x})=\Normal(\boldsymbol\mu,\boldsymbol\Sigma)$, by minimizing
$\KL{p}{q}$ with respect to $\boldsymbol\mu$ and $\boldsymbol\Sigma$. Which expressions minimize the
KL-divergence? ($\E$ is expectation, $\Cov$ the covariance, $\mathbf{I}$ the identity matrix.)
\begin{itemize}[leftmargin=1.8em]
  \choice $\boldsymbol\mu=\E(\mathbf{x})$ and $\boldsymbol\Sigma=\Cov(\mathbf{x})$
  \choice $\boldsymbol\mu=\E(\mathbf{x})$ and $\boldsymbol\Sigma=\mathbf{I}$
  \choice $\boldsymbol\mu=\E(\mathbf{x}\T\mathbf{x})$ and $\boldsymbol\Sigma=\E(\mathbf{x}\T)$
  \choice $\boldsymbol\mu=\E(\mathbf{x}\T\mathbf{x})$ and $\boldsymbol\Sigma=\mathbf{I}$
  \choice $\boldsymbol\mu=\E(\mathbf{x})$ and $\boldsymbol\Sigma=\E(\mathbf{x}\mathbf{x}\T)$
\end{itemize}
\closequestion
\begin{answer}{A}
  \yes{A} Minimizing $\KL{p}{q}$ over a Gaussian \emph{matches the first two moments}: $\boldsymbol\mu=\E(\mathbf{x})$, $\boldsymbol\Sigma=\Cov(\mathbf{x})$.
  \no{B} Forcing $\boldsymbol\Sigma=\mathbf{I}$ throws away the true covariance.
  \no{C} $\boldsymbol\mu=\E(\mathbf{x}\T\mathbf{x})$ is a scalar --- dimensionally nonsense.
  \no{D} Same dimensional nonsense for $\boldsymbol\mu$.
  \no{E} Near-miss: $\E(\mathbf{x}\mathbf{x}\T)$ is the \emph{second moment}, which over-estimates the covariance because it omits the $-\E(\mathbf{x})\E(\mathbf{x})\T$ correction.
\end{answer}


% ==========================================================
\section{GANs and Wasserstein GANs}
% ==========================================================

\questiontitle{GAN Objective Function}
Using the lecture notation, which of the following is the correct formulation of the standard
objective of Generative Adversarial Networks (GANs)?
\begin{itemize}[leftmargin=1.8em]
  \choice $\min_G\max_D \loss(D,G)=\E_{\mathbf{x}\sim p_r}[\log D(\mathbf{x})]+\E_{\mathbf{z}\sim p_z}[\log(1-D(\mathbf{z}))]$
  \choice $\min_G\max_D \loss(D,G)=\E_{\mathbf{x}\sim p_r}[\log D(\mathbf{x})]+\E_{\mathbf{x}\sim p_r}[\log(1-D(G(\mathbf{x})))]$
  \choice $\min_G\max_D \loss(D,G)=\E_{\mathbf{x}\sim p_r}[\log G(\mathbf{x})]+\E_{\mathbf{z}\sim p_z}[\log(1-G(D(\mathbf{z})))]$
  \choice $\min_G\max_D \loss(D,G)=\E_{\mathbf{x}\sim p_r}[\log D(\mathbf{x})]+\E_{\mathbf{z}\sim p_z}[\log(1-D(G(\mathbf{z})))]$
  \choice $\min_G\max_D \loss(D,G)=\E_{\mathbf{x}\sim p_r}[\log D(\mathbf{x})]+\E_{\mathbf{z}\sim p_z}[D(G(\mathbf{z}))]$
  \choice $\min_G\max_D \loss(D,G)=\E_{\mathbf{x}\sim p_r}[\log D(\mathbf{x})]+\E_{\mathbf{z}\sim p_z}[\log(D(G(\mathbf{z})))]$
\end{itemize}
\closequestion
\begin{answer}{D}
  \no{A} The fake term must pass through $G$: $D(G(\mathbf{z}))$, not $D(\mathbf{z})$.
  \no{B} The second expectation must be over $\mathbf{z}\sim p_z$, not over real data $\mathbf{x}\sim p_r$.
  \no{C} Roles of $D$ and $G$ are swapped and the wrong functions are logged --- nonsense.
  \yes{D} Correct value function: $\E_{p_r}[\log D(\mathbf{x})]+\E_{p_z}[\log(1-D(G(\mathbf{z})))]$.
  \no{E} Missing the $\log(1-\cdot)$ on the fake term.
  \no{F} Wrong form: should be $\log(1-D(G(\mathbf{z})))$, not $\log(D(G(\mathbf{z})))$.
\end{answer}

\questiontitle{Properties of a Trained GAN}
Consider a GAN that successfully produces images of apples. Which statements are correct?
(Select one or more.)
\begin{itemize}[leftmargin=1.8em]
  \choice The discriminator can be used to classify images as apple vs.\ non-apple.
  \choice The generator aims to learn the distribution of apple images.
  \choice The generator can produce unseen images of apples.
  \choice After training the GAN, the discriminator loss eventually reaches a constant value.
\end{itemize}
\closequestion
\begin{answer}{B, C}
  \no{A} The discriminator only learns \emph{real-vs-fake} apples, not apple-vs-non-apple.
  \yes{B} The generator's goal is exactly to model the distribution of apple images.
  \yes{C} A good generator generalizes and produces novel, unseen apples.
  \no{D} Not guaranteed: at the ideal equilibrium $D=0.5$ the loss would be constant, but in general GAN training need not converge there.
\end{answer}

\questiontitle{Divergences and Distances}
Which of the following statements about divergences and distances are correct? (Select one or more.)
\begin{itemize}[leftmargin=1.8em]
  \choice JS and KL divergences show problems for distributions with disjoint supports --- the Wasserstein
        distance overcomes this.
  \choice A map $D:\mathcal{S}\times\mathcal{S}\to\R$ is a divergence iff it is non-negative and symmetric.
  \choice The Wasserstein-1 distance with Euclidean cost is a distance, not a divergence.
  \choice The Wasserstein distance is defined as a minimization over joint distributions (couplings),
        but this is infeasible in practice.
  \choice The Kantorovich--Rubinstein theorem lets us write the Wasserstein distance as a maximization
        over Lipschitz functions, which makes it tractable.
\end{itemize}
\closequestion
\begin{answer}{A, C, D, E}
  \yes{A} KL / JS blow up or stay constant for disjoint supports; Wasserstein still gives a finite, smooth value.
  \no{B} A divergence need \emph{not} be symmetric --- non-negative $+$ symmetric (with the triangle inequality) describes a \emph{distance/metric}, not a divergence.
  \yes{C} $W_1$ with Euclidean cost satisfies the metric axioms (symmetry, triangle inequality), so it is a genuine distance.
  \yes{D} It is an infimum over couplings (joint distributions) --- intractable to compute directly.
  \yes{E} Kantorovich--Rubinstein duality rewrites it as a supremum over $1$-Lipschitz functions --- exactly what WGAN exploits.
\end{answer}

\questiontitle{KL, JSD and Wasserstein vs.\ Separation}
Let $p_1=\Normal(0,1)$ and $p_2=\Normal(x,1)$. As $x$ increases, you plot the
KL-divergence $\KL{p_1}{p_2}$, the Jensen--Shannon divergence $\mathrm{JSD}(p_1\|p_2)$, and the
Wasserstein-2 distance $\mathrm{WD}(p_1\|p_2)$ against $x$. One curve grows quadratically without
bound, one grows linearly, and one saturates at a finite value. Match each measure to its qualitative
behaviour.
\examfigure[width=0.86\linewidth]{figures/exam-kl-jsd-wasserstein-separation.png}
\begin{itemize}[leftmargin=1.8em]
  \item KL-divergence
  \item Wasserstein-2 distance
  \item Jensen--Shannon divergence
\end{itemize}
\closequestion
\begin{answerprose}{KL quadratic, $W_2$ linear, JSD saturates}
\begin{itemize}[leftmargin=1.7em,itemsep=1.5pt,topsep=2pt]
  \item[\textcolor{ok}{$\to$}] \textbf{KL-divergence:} grows \emph{quadratically} without bound --- for equal-variance Gaussians $\KL{p_1}{p_2}=x^2/2$.
  \item[\textcolor{ok}{$\to$}] \textbf{Wasserstein-2:} grows \emph{linearly} --- $\mathrm{WD}_2=|x|$.
  \item[\textcolor{ok}{$\to$}] \textbf{Jensen--Shannon:} \emph{saturates} at a finite value (bounded by $\log 2$) once the supports barely overlap.
\end{itemize}
This is why JS/KL give vanishing gradients for far-apart distributions while Wasserstein stays informative.
\end{answerprose}


% ==========================================================
\section{Metrics, Normalizing Flows \& Density Models}
% ==========================================================

\questiontitle{Inception Score and FID}
Tick the correct statements about the Inception Score (IS) and the Fréchet Inception Distance (FID).
(Select one or more.)
\begin{itemize}[leftmargin=1.8em]
  \choice FID can detect mode collapse, because mode collapse influences the covariance matrix of the
        generated data.
  \choice IS and FID are metrics that assess the quality of GAN-generated images.
  \choice Both IS and FID rely on features of the input images produced by AlexNet.
  \choice IS and FID can be used as GAN objectives that lead to faster training.
  \choice The computation of the FID is based on an explicit formula for the Wasserstein-2 distance
        between two Gaussian distributions.
\end{itemize}
\closequestion
\begin{answer}{A, B, E}
  \yes{A} FID compares means \emph{and covariances}; mode collapse shrinks the covariance, so FID picks it up.
  \yes{B} Both are standard metrics for the quality of GAN-generated images.
  \no{C} They use \emph{Inception} (v3) features, not AlexNet.
  \no{D} They are \emph{evaluation} metrics, not training objectives.
  \yes{E} FID is exactly the closed-form Wasserstein-2 (Fréchet) distance between two Gaussians fitted to the features.
\end{answer}

\questiontitle{Normalizing Flows}
Which of the following statements about normalizing flows are correct? (Select one or more.)
\begin{itemize}[leftmargin=1.8em]
  \choice It is impossible to integrate a multi-layer perceptron as a transformation in a normalizing
        flow because MLPs are not invertible.
  \choice All multi-layer perceptrons with the same input and output dimension can be considered a
        normalizing flow.
  \choice Normalizing flows cannot contain functions with a lower-triangular Jacobian because they are
        not invertible.
  \choice Some normalizing flows use functions with a lower-triangular Jacobian to facilitate the
        computation of the determinant.
\end{itemize}
\closequestion
\begin{answer}{D}
  \no{A} ``Impossible'' is too strong --- one \emph{can} design invertible MLP-based transforms; ordinary MLPs simply aren't invertible by default.
  \no{B} Matching input/output dimensions does not make a network invertible / a valid flow.
  \no{C} Lower-triangular-Jacobian maps can absolutely be invertible --- that is the whole point of autoregressive flows.
  \yes{D} A triangular Jacobian makes the determinant the product of its diagonal --- cheap to compute.
\end{answer}

\questiontitle{Flow-Based Models}
Which of the following statements about flow-based models are correct? (Select one or more.)
\begin{itemize}[leftmargin=1.8em]
  \choice Theoretically, flow-based models can represent any distribution $p_x(\mathbf{x})$ if one starts
        from a sufficiently complex distribution; a simple base distribution $p_u(\mathbf{u})$ such as the
        uniform on the cube $(0,1)^D$ would not be sufficient.
  \choice Fitting flow-based models is often done with the KL-divergence.
  \choice Computing the Jacobi determinant of a neural network with $D$ inputs and $D$ outputs is in
        general of $\mathcal{O}(D^3)$.
  \choice For autoregressive flows the transformation $g_\phi$ is constructed so that its Jacobi
        determinant always has value 1.
\end{itemize}
\closequestion
\begin{answer}{B, C}
  \no{A} A \emph{simple} base (even uniform on $(0,1)^D$ or a standard Gaussian) is sufficient --- the flexible transform does the work.
  \yes{B} Flows are fit by maximum likelihood, i.e.\ minimizing a KL-divergence to the data.
  \yes{C} A general dense Jacobian determinant costs $\mathcal{O}(D^3)$.
  \no{D} An autoregressive flow's Jacobian is triangular, but its determinant (product of the diagonal) is \emph{not} fixed at $1$.
\end{answer}

\questiontitle{Density under a Linear Transform}
\repeatnote\\
Assume a two-dimensional vector $\mathbf{u}\sim p_u(\mathbf{u})$. Let
$T=\begin{pmatrix}2&0\\0&\tfrac12\end{pmatrix}$ and $\mathbf{x}=T\cdot\mathbf{u}$. Mark the correct
formula for the distribution $p_x(\mathbf{x})$.
{\footnotesize
\begin{itemize}[leftmargin=1.8em]
  \choice $p_u\!\left(\begin{pmatrix}\tfrac12&0\\0&2\end{pmatrix}\mathbf{x}\right)\cdot\tfrac12$
  \choice $p_u\!\left(\begin{pmatrix}\tfrac12&0\\0&2\end{pmatrix}\mathbf{x}\right)$
  \choice $p_u\!\left(\begin{pmatrix}\tfrac12&0\\0&1\end{pmatrix}\mathbf{x}\right)\cdot\tfrac12$
  \choice $p_u\!\left(\begin{pmatrix}\tfrac12&0\\0&1\end{pmatrix}\mathbf{x}\right)$
  \choice $p_u\!\left(\begin{pmatrix}2&0\\0&1\end{pmatrix}\mathbf{x}\right)\cdot 2$
  \choice $p_u\!\left(\begin{pmatrix}2&0\\0&1\end{pmatrix}\mathbf{x}\right)\cdot\tfrac12$
\end{itemize}}
\closequestion
\begin{answer}{B}
  \yes{B} Change of variables: $p_x(\mathbf{x})=\frac{1}{|\det T|}\,p_u(T^{-1}\mathbf{x})$ with $T^{-1}=\begin{psmallmatrix}\tfrac12&0\\0&2\end{psmallmatrix}$ and $\det T=2\cdot\tfrac12=1$, so the prefactor is $1$.
  \no{A} Correct $T^{-1}$ but a spurious $\tfrac12$ factor ($\det T=1$, not $2$).
  \no{C--F} Wrong inverse matrix and/or wrong determinant prefactor --- guess-catchers.
\end{answer}

\questiontitle{Density under a Linear Transform (Non-Diagonal)}
Let $\mathbf{u}\sim p_u(\mathbf{u})$ be a two-dimensional random vector and define
$\mathbf{x}=T\mathbf{u}$ with the non-diagonal matrix
$T=\begin{pmatrix}2&-2\\0&1\end{pmatrix}$. Derive the density $p_x(\mathbf{x})$ in terms of $p_u$,
stating $\det T$, $T^{-1}$, and the final formula.
\closequestion
\begin{answerprose}{$p_x(\mathbf{x})=\tfrac12\,p_u(T^{-1}\mathbf{x})$, $\det T=2$}
$\det T = 2\cdot1-(-2)\cdot0 = 2$, hence
$T^{-1}=\tfrac12\begin{pmatrix}1&2\\0&2\end{pmatrix}=\begin{pmatrix}\tfrac12&1\\0&1\end{pmatrix}$.
Change of variables gives
$p_x(\mathbf{x})=\frac{1}{|\det T|}\,p_u(T^{-1}\mathbf{x})=\tfrac12\,p_u\!\left(\begin{pmatrix}\tfrac12&1\\0&1\end{pmatrix}\mathbf{x}\right)$.
Unlike the diagonal case, $T$ is triangular but \emph{not} diagonal, so $T^{-1}$ carries an
off-diagonal entry and the prefactor is $\tfrac12$ (not $1$).
\end{answerprose}

\questiontitle{Autoregressive Sequence Probability}
An autoregressive model $p_w(\mathbf{x})$ models sequences $(x_1,x_2,x_3)$ where each element is
$A$, $B$, or $C$:
\[
  p_w(x_1)=\tfrac13,\qquad
  p_w(x_2\given x_1)=\begin{cases}\tfrac12 & x_2=x_1\\[2pt]\tfrac14 & x_2\neq x_1\end{cases},\qquad
  p_w(x_3\given x_2,x_1)=\begin{cases}\tfrac12 & x_3=x_2\\[2pt]\tfrac14 & x_3\neq x_2\end{cases}.
\]
Compute $p_w(\mathbf{x})$ for $\mathbf{x}=(A,A,C)$, rounded to three decimals.
\closequestion
\begin{answerprose}{$1/24 \approx 0.042$}
$p_w(A,A,C)=\underbrace{p_w(x_1{=}A)}_{1/3}\cdot\underbrace{p_w(x_2{=}A\mid A)}_{1/2\ (x_2=x_1)}\cdot\underbrace{p_w(x_3{=}C\mid A)}_{1/4\ (x_3\neq x_2)}=\tfrac13\cdot\tfrac12\cdot\tfrac14=\tfrac{1}{24}\approx 0.042.$
\end{answerprose}


% ==========================================================
\section{Bayesian Deep Learning}
% ==========================================================

\questiontitle{Deep Ensembles and Marginalization}
The main difference between the frequentist and Bayesian approach is the marginalization over
parameters when defining the probability of a class label given input $\mathbf{x}$ and dataset
$\mathcal{D}$:
\[
  p(\mathbf{y}^{\mathrm{test}}\given\mathbf{x}^{\mathrm{test}},\mathcal{D})
  = \E_{\mathbf{w}\sim p(\mathbf{w}\given\mathcal{D})}\!\left[p(\mathbf{y}^{\mathrm{test}}\given\mathbf{w},\mathbf{x}^{\mathrm{test}})\right].
\]
Deep Ensembles approximate this full expectation with (weighted) averages. Which of the following
are valid approximations? (Select one or more.)
\begin{itemize}[leftmargin=1.8em]
  \choice $\tfrac1M\sum_{m=1}^{M} p(\mathbf{y}^{\mathrm{test}}\given\mathbf{w}_m,\mathbf{x}^{\mathrm{test}})$,
        where $\mathbf{w}_m$ are the weights of the $m$-th network of the ensemble.
  \choice $p(\mathbf{y}^{\mathrm{test}}\given\mathbf{w},\mathbf{x}^{\mathrm{test}})$ for a single trained network.
  \choice $\tfrac1M\sum_{m=1}^{M}\alpha_m\, p(\mathbf{y}^{\mathrm{test}}\given\mathbf{w}_m,\mathbf{x}^{\mathrm{test}})$
        with $\alpha_m=\softmax(\mathbf{z})_m$ for some random vector $\mathbf{z}$.
  \choice $\tfrac1M\sum_{m=1}^{M} p(\mathbf{y}^{\mathrm{test}}_m\given\mathbf{w},\mathbf{x}^{\mathrm{test}}_m)$,
        where $m$ denotes the $m$-th equally-sized fraction of the test data.
\end{itemize}
\closequestion
\begin{answer}{A, C}
  \yes{A} A uniform average over the $M$ ensemble members approximates the posterior expectation.
  \no{B} A single network is just one sample --- no marginalization over the posterior.
  \yes{C} A \emph{weighted} average (weights from a softmax, summing to $1$) is also a valid approximation of the expectation.
  \no{D} This splits the \emph{test data} into fractions --- that is not averaging over parameters at all.
\end{answer}


% ==========================================================
\section{Graph Neural Networks}
% ==========================================================

\questiontitle{Distinguishing Two Small Graphs}
Consider two graphs: $\mathcal{G}_1$ is a path of three nodes (a coloured centre node joined to two
leaves), and $\mathcal{G}_2$ is a star with a centre node joined to three leaves. Mark the correct
statements. (Select one or more.)
\examfigure[width=0.62\linewidth]{figures/exam-gnn-small-graphs.png}
\begin{itemize}[leftmargin=1.8em]
  \choice $\mathcal{G}_1$ and $\mathcal{G}_2$ can be distinguished by the Weisfeiler--Lehman
        (colour refinement) test.
  \choice All message-passing neural networks (MPNNs) can distinguish $\mathcal{G}_1$ and $\mathcal{G}_2$.
  \choice $\mathcal{G}_1$ and $\mathcal{G}_2$ are isomorphic graphs.
\end{itemize}
\closequestion
\begin{answer}{A}
  \yes{A} They have different structure (different node counts / degree sequences), so colour refinement separates them.
  \no{B} ``All MPNNs'' is too strong --- expressivity depends on the aggregation; a weak MPNN need not distinguish arbitrary non-isomorphic graphs.
  \no{C} They are not isomorphic (different structure).
\end{answer}

\questiontitle{Weisfeiler--Lehman Necessary Condition}
Two graphs $A$ and $B$ are given. Do they fulfil the necessary condition for graph isomorphism
given by the Weisfeiler--Lehman test (i.e.\ do they produce the same stable colour histogram)?
\begin{itemize}[leftmargin=1.8em]
  \choice True
  \choice False
\end{itemize}
\closequestion
\begin{answer}{True}
  \yes{True} The two graphs produce the same stable WL colour histogram, so the \emph{necessary} condition for isomorphism is satisfied. (Necessary, not sufficient --- equal histograms don't prove isomorphism.)
  \no{False} ---
\end{answer}

\questiontitle{Matching Graphs to Adjacency Matrices}
Four graphs on six nodes are shown: $A$ is a plain $6$-cycle (hexagon), $B$ is two triangles,
$C$ is another $6$-cycle drawn differently, and $D$ is a triangle with a three-edge tail attached.
For each adjacency matrix, name the matching graph ($A$, $B$, $C$, $D$, or none).
\examfigure[width=0.94\linewidth]{figures/exam-gnn-candidate-graphs.png}
\begin{itemize}[leftmargin=1.8em]
  \item All six row-sums equal $2$ and the edges form a single closed cycle.
  \item Two row-sums equal $1$ and the rest equal $2$ (a single open path with two endpoints).
  \item All six row-sums equal $2$ but the edges split into two separate $3$-cycles.
  \item Four row-sums equal $2$ and two equal $3$ (a cycle plus one chord).
\end{itemize}
\closequestion
\begin{answerprose}{$A/C$ \ / \ none \ / \ $B$ \ / \ none}
Read off each matrix by its row-sums (node degrees) and connectivity:
\begin{itemize}[leftmargin=1.7em,itemsep=1.5pt,topsep=2pt]
  \item[\textcolor{ok}{$\to$}] All degrees $2$, one closed cycle $\Rightarrow$ a $6$-cycle: \textbf{$A$} (and the differently-drawn cycle is \textbf{$C$}).
  \item[\textcolor{ok}{$\to$}] Two degree-$1$ endpoints (single open path) $\Rightarrow$ \textbf{none of the above}. This is not \textbf{$B$}: $B$ has two separate $3$-cycles, so every node in $B$ has degree $2$.
  \item[\textcolor{ok}{$\to$}] All degrees $2$ but split into two $3$-cycles $\Rightarrow$ \textbf{$B$} (two triangles).
  \item[\textcolor{ok}{$\to$}] Two degree-$3$ nodes (cycle $+$ chord) $\Rightarrow$ \textbf{none of the above}. The shown \textbf{$D$} has one degree-$1$ endpoint and one degree-$3$ attachment node, so it is a triangle with a tail rather than a chorded $6$-cycle.
\end{itemize}
\end{answerprose}

\questiontitle{Group Action in Graph NNs}
Which of the following is a common group action used in graph neural networks?
\begin{itemize}[leftmargin=1.8em]
  \choice Rotation
  \choice Scaling
  \choice Translation (permutation / relabelling of nodes)
  \choice Cropping
  \choice Clustering
\end{itemize}
\closequestion
\begin{answer}{C}
  \no{A} Image augmentation, not the GNN symmetry.
  \no{B} Image augmentation, not relevant here.
  \yes{C} GNNs are built to be invariant / equivariant to \emph{permutations} (relabelling) of the nodes.
  \no{D} Image augmentation.
  \no{E} A downstream task, not a group action.
\end{answer}

\questiontitle{Equivariance of an Attention-Pooling Layer}
A layer implements $\mathbf{h}^{(l+1)}=M\,\underbrace{\softmax(\beta\,M\T\mathbf{h}^{(l)})}_{=:\,\mathbf{p}}$,
where $\mathbf{h}^{(l)}\in\R^{d}$, $\beta>0$ is a hyperparameter, and $M\in\R^{d\times J}$ holds $J$
pattern columns. Which statements about this layer are true? (Select one or more.)
\begin{itemize}[leftmargin=1.8em]
  \choice $\mathbf{h}^{(l+1)}$ is equivariant to permutations of the columns of $M$.
  \choice $\mathbf{p}$ is invariant to permutations of the columns of $M$.
  \choice $\mathbf{p}$ is equivariant to permutations of the columns of $M$.
  \choice $\mathbf{h}^{(l+1)}$ is invariant to permutations of the columns of $M$.
  \choice $\mathbf{p}$ is equivariant to permutations of the rows of $M$.
  \choice This layer cannot be used in deep learning because the gradient
        $\partial\mathbf{h}^{(l+1)}/\partial\mathbf{h}^{(l)}$ cannot be backpropagated.
\end{itemize}
\closequestion
\begin{answer}{C, D}
  \no{A} $\mathbf{h}^{(l+1)}$ is \emph{invariant}, not equivariant, to column permutations of $M$.
  \no{B} $\mathbf{p}$'s entries permute under column permutations of $M$, so it is not invariant.
  \yes{C} Permuting the columns of $M$ permutes the rows of $M\T$, hence permutes the entries of $\mathbf{p}$ --- $\mathbf{p}$ is \emph{equivariant}.
  \yes{D} $\mathbf{h}^{(l+1)}=M\mathbf{p}$ is unchanged: the column reorder of $M$ cancels the row reorder of $\mathbf{p}$ --- so it is \emph{invariant}.
  \no{E} Row permutations act on the feature dimension $d$, not on the $J$ pattern entries of $\mathbf{p}$ --- this does not give equivariance of $\mathbf{p}$.
  \no{F} The layer (softmax $+$ matrix products) is differentiable --- gradients backpropagate fine.
\end{answer}


% ==========================================================
\section{Theory \& Foundations}
% ==========================================================

\questiontitle{Universal Approximator Theorem}
Mark the correct statements about the Universal Approximator Theorem (UAT) for neural networks.
(Select one or more.)
\begin{itemize}[leftmargin=1.8em]
  \choice For a real-valued, continuous function $h$ on $[-\pi,\pi]^2$, the existence of an approximating
        network can be derived with the help of Fourier series.
  \choice The UAT is a theoretical statement about the expressive power of neural networks.
  \choice The UAT gives a useful algorithm to compute the optimal weights of a network that should
        approximate a given continuous function.
  \choice The algorithm related to the UAT is easy to understand and implement, but very costly and
        therefore hardly used in practice.
  \choice The UAT shows that any function can be approximated by any given network with arbitrary
        precision; there are no restrictions on the function or on the network structure.
\end{itemize}
\closequestion
\begin{answer}{A, B}
  \yes{A} One proof route builds the approximation from Fourier series on the compact domain.
  \yes{B} The UAT is purely an existence / expressivity statement.
  \no{C} It is \emph{non-constructive} --- it does not tell you how to find the weights.
  \no{D} There is no practical ``UAT algorithm'' in use --- nonsense.
  \no{E} There are conditions: the target must be continuous on a compact set and the network needs at least one (sufficiently wide) hidden layer --- ``any function, no restrictions'' is false.
\end{answer}

\questiontitle{Newton's Method vs.\ Gradient Direction}
On a quadratic error surface (elliptical level curves), two arrows are drawn from a point: one points
straight at the minimum at the centre of the ellipses, the other points perpendicular to the local
level curve. Label each arrow as \emph{Gradient} or \emph{Newton's Method}.
\examfigure[width=0.9\linewidth]{figures/exam-newton-gradient-directions.png}
\begin{itemize}[leftmargin=1.8em]
  \item Arrow pointing directly to the centre (minimum)
  \item Arrow perpendicular to the level curve (steepest descent)
\end{itemize}
\closequestion
\begin{answerprose}{centre $=$ Newton, perpendicular $=$ Gradient}
\begin{itemize}[leftmargin=1.7em,itemsep=1.5pt,topsep=2pt]
  \item[\textcolor{ok}{$\to$}] \textbf{Arrow to the centre $=$ Newton's Method.} It pre-multiplies the gradient by the inverse Hessian, so on a quadratic it points straight at the minimum in one step.
  \item[\textcolor{ok}{$\to$}] \textbf{Arrow perpendicular to the level curve $=$ Gradient.} Steepest descent is orthogonal to the contour, which generally does \emph{not} aim at the minimum on elongated ellipses.
\end{itemize}
\end{answerprose}


\end{document}
